% --- Archon physics formalization source begin ---
% archon:physics
% archon:covers PhyXMiniProblems/problem_phyx_mini_0631.lean
% archon:source-report reports/phyx_mini/problem_phyx_mini_0631.source.json

\chapter{Physics problem phyx\_mini\_0631}
\label{ch:PhyXMiniProblems_problem_phyx_mini_0631}

\paragraph{Problem source.}
\#\# Physical scenario

A colorless, odorless gas in a tank needs to be identified. The gas is nitrogen (\$N\_2\$), carbon monoxide (\$CO\$), or nitric oxide (\$NO\$). White light is passed through a sample to obtain an absorption spectrum. Among others, the following wavelengths present dark bands on the spectrum: \$\textbackslash{}lambda = 4.2680\textbackslash{} \textbackslash{}mu m, 4.2753\textbackslash{} \textbackslash{}mu m, 4.2826\textbackslash{} \textbackslash{}mu m, 4.2972\textbackslash{} \textbackslash{}mu m\$, and \$4.3046\textbackslash{} \textbackslash{}mu m\$. These represent transitions from the lowest rotational states.

\#\# Question

Which one of the absorbed energies corresponds to a transition to an \$l=0\$ state?

\#\# Answer choices

- A: A: 0.288528eV

- B: B: 0.298528eV

- C: C: 0.281258eV

- D: D: 0.282145eV

\#\# Figure caption (auxiliary; use the image as primary evidence)

The image shows three molecular diagrams, each depicting a pair of atoms connected by a line representing a chemical bond. 

1. **Top Diagram**:
   - Atoms: Two nitrogen atoms (N) are connected.
   - Atomic Mass Numbers: Each nitrogen atom is labeled with the number 14.
   - Bond Length: The bond between them is labeled as 110 pm.

2. **Middle Diagram**:
   - Atoms: One carbon (C) atom and one oxygen (O) atom are connected.
   - Atomic Mass Numbers: The carbon atom is labeled with the number 12, and the oxygen atom with the number 16.
   - Bond Length: The bond is labeled as 118 pm.

3. **Bottom Diagram**:
   - Atoms: One nitrogen (N) atom and one oxygen (O) atom are connected.
   - Atomic Mass Numbers: The nitrogen atom is labeled with the number 14, and the oxygen atom with the number 16.
   - Bond Length: The bond is labeled as 115 pm.

The atoms are represented by colored spheres (blue for nitrogen, black for carbon, and red for oxygen), and the numbers are presented within the spheres. The bond lengths are displayed as text above the bonds.

\#\# Dataset metadata

Quantum Phenomena; Physical Model Grounding Reasoning, Multi-Formula Reasoning

\paragraph{Recorded answer/context.}
A: A: 0.288528eV

\paragraph{Figure/image path.}
/root/proposal\_for\_physic/science-mango/phyx\_mini\_run/phyx\_data/test\_image/631.png

\paragraph{Formalization target.}
create a compiling Lean file with sorry bodies at `PhyXMiniProblems/problem\_phyx\_mini\_0631.lean`. The Lean declarations must preserve the physical quantities, dimensions or dimensional roles, figure labels, governing-law hypotheses, and final relation expressed by this problem.
Use Mathlib/Physlib names found through LeanExplore where available. If a physics API is missing, introduce faithful local abstractions rather than scalar placeholder aliases.

\begin{theorem}[Physics formalization target]
\label{thm:physics:phyx_mini_0631:target}
\lean{PhyXMiniProblems.ProblemPhyXMini0631.problem_phyx_mini_0631}
\uses{def:physics:phyx-mini-0631:phyxminiproblems-problemphyxmini0631-diatomicabsorptionsetup, def:physics:phyx-mini-0631:phyxminiproblems-problemphyxmini0631-matchesgasabsorptionscenario, def:physics:phyx-mini-0631:phyxminiproblems-problemphyxmini0631-matchesprimarydiatomiccandidatesfigure, def:physics:phyx-mini-0631:phyxminiproblems-problemphyxmini0631-matchesgivenabsorptionwavelengths, def:physics:phyx-mini-0631:phyxminiproblems-problemphyxmini0631-usestextbookreferencedata, def:physics:phyx-mini-0631:phyxminiproblems-problemphyxmini0631-hasphysicaldiatomicabsorptionparameters, def:physics:phyx-mini-0631:phyxminiproblems-problemphyxmini0631-satisfiesdiatomicmomentofinertialaws, def:physics:phyx-mini-0631:phyxminiproblems-problemphyxmini0631-satisfiesrigidrotorabsorptionlaws, def:physics:phyx-mini-0631:phyxminiproblems-problemphyxmini0631-agreeswithdisplayedenergy, def:physics:phyx-mini-0631:phyxminiproblems-problemphyxmini0631-isuniqueclosestenergychoice}
The assigned autoformalize agent should translate this physics problem into Lean declarations in the covered file, with theorem and lemma proof bodies written as `by sorry`.
\end{theorem}
\begin{proof}
This is an autoformalization task, not a proof task. Produce statements that can later be proved by the physics prover without weakening the physical model.
\end{proof}

% --- Archon declaration topology begin ---
\section{Declaration topology}
\begin{definition}[Length Quantity]
  \label{def:physics:phyx-mini-0631:phyxminiproblems-problemphyxmini0631-lengthquantity}
  \lean{PhyXMiniProblems.ProblemPhyXMini0631.LengthQuantity}
  A nonnegative, unit-independent physical length
\end{definition}
\begin{proof}
  This is the declared model definition of Length Quantity.
\end{proof}
\begin{definition}[Mass Quantity]
  \label{def:physics:phyx-mini-0631:phyxminiproblems-problemphyxmini0631-massquantity}
  \lean{PhyXMiniProblems.ProblemPhyXMini0631.MassQuantity}
  A nonnegative, unit-independent physical mass
\end{definition}
\begin{proof}
  This is the declared model definition of Mass Quantity.
\end{proof}
\begin{definition}[moment Of Inertia Dimension]
  \label{def:physics:phyx-mini-0631:phyxminiproblems-problemphyxmini0631-momentofinertiadimension}
  \lean{PhyXMiniProblems.ProblemPhyXMini0631.momentOfInertiaDimension}
  The dimension of moment of inertia, mass times length squared
\end{definition}
\begin{proof}
  This is the declared model definition of moment Of Inertia Dimension.
\end{proof}
\begin{definition}[Moment Of Inertia Quantity]
  \label{def:physics:phyx-mini-0631:phyxminiproblems-problemphyxmini0631-momentofinertiaquantity}
  \lean{PhyXMiniProblems.ProblemPhyXMini0631.MomentOfInertiaQuantity}
  \uses{def:physics:phyx-mini-0631:phyxminiproblems-problemphyxmini0631-momentofinertiadimension}
  A nonnegative, unit-independent moment of inertia
\end{definition}
\begin{proof}
  This is the declared model definition of Moment Of Inertia Quantity.
\end{proof}
\begin{definition}[action Dimension]
  \label{def:physics:phyx-mini-0631:phyxminiproblems-problemphyxmini0631-actiondimension}
  \lean{PhyXMiniProblems.ProblemPhyXMini0631.actionDimension}
  The dimension of action, equivalently energy times time
\end{definition}
\begin{proof}
  This is the declared model definition of action Dimension.
\end{proof}
\begin{definition}[Action Quantity]
  \label{def:physics:phyx-mini-0631:phyxminiproblems-problemphyxmini0631-actionquantity}
  \lean{PhyXMiniProblems.ProblemPhyXMini0631.ActionQuantity}
  \uses{def:physics:phyx-mini-0631:phyxminiproblems-problemphyxmini0631-actiondimension}
  A nonnegative, unit-independent action such as reduced Planck action
\end{definition}
\begin{proof}
  This is the declared model definition of Action Quantity.
\end{proof}
\begin{definition}[Speed Quantity]
  \label{def:physics:phyx-mini-0631:phyxminiproblems-problemphyxmini0631-speedquantity}
  \lean{PhyXMiniProblems.ProblemPhyXMini0631.SpeedQuantity}
  A signed-carrier speed quantity, matching the exact type of Physlib's DimSpeed.speedOfLight
\end{definition}
\begin{proof}
  This is the declared model definition of Speed Quantity.
\end{proof}
\begin{definition}[length Readout]
  \label{def:physics:phyx-mini-0631:phyxminiproblems-problemphyxmini0631-lengthreadout}
  \lean{PhyXMiniProblems.ProblemPhyXMini0631.lengthReadout}
  \uses{def:physics:phyx-mini-0631:phyxminiproblems-problemphyxmini0631-lengthquantity}
  Read a physical length as a real number in a selected length unit
\end{definition}
\begin{proof}
  This is the declared model definition of length Readout.
\end{proof}
\begin{definition}[length In Meters]
  \label{def:physics:phyx-mini-0631:phyxminiproblems-problemphyxmini0631-lengthinmeters}
  \lean{PhyXMiniProblems.ProblemPhyXMini0631.lengthInMeters}
  \uses{def:physics:phyx-mini-0631:phyxminiproblems-problemphyxmini0631-lengthquantity, def:physics:phyx-mini-0631:phyxminiproblems-problemphyxmini0631-lengthreadout}
  Metre readout of a physical length
\end{definition}
\begin{proof}
  This is the declared model definition of length In Meters.
\end{proof}
\begin{definition}[length In Picometers]
  \label{def:physics:phyx-mini-0631:phyxminiproblems-problemphyxmini0631-lengthinpicometers}
  \lean{PhyXMiniProblems.ProblemPhyXMini0631.lengthInPicometers}
  \uses{def:physics:phyx-mini-0631:phyxminiproblems-problemphyxmini0631-lengthquantity, def:physics:phyx-mini-0631:phyxminiproblems-problemphyxmini0631-lengthreadout}
  Picometre readout of a physical length
\end{definition}
\begin{proof}
  This is the declared model definition of length In Picometers.
\end{proof}
\begin{definition}[length In Micrometers]
  \label{def:physics:phyx-mini-0631:phyxminiproblems-problemphyxmini0631-lengthinmicrometers}
  \lean{PhyXMiniProblems.ProblemPhyXMini0631.lengthInMicrometers}
  \uses{def:physics:phyx-mini-0631:phyxminiproblems-problemphyxmini0631-lengthquantity, def:physics:phyx-mini-0631:phyxminiproblems-problemphyxmini0631-lengthreadout}
  Micrometre readout of a physical length
\end{definition}
\begin{proof}
  This is the declared model definition of length In Micrometers.
\end{proof}
\begin{definition}[mass In Kilograms]
  \label{def:physics:phyx-mini-0631:phyxminiproblems-problemphyxmini0631-massinkilograms}
  \lean{PhyXMiniProblems.ProblemPhyXMini0631.massInKilograms}
  \uses{def:physics:phyx-mini-0631:phyxminiproblems-problemphyxmini0631-massquantity}
  Kilogram readout of a physical mass
\end{definition}
\begin{proof}
  This is the declared model definition of mass In Kilograms.
\end{proof}
\begin{definition}[mass In Atomic Mass Units]
  \label{def:physics:phyx-mini-0631:phyxminiproblems-problemphyxmini0631-massinatomicmassunits}
  \lean{PhyXMiniProblems.ProblemPhyXMini0631.massInAtomicMassUnits}
  \uses{def:physics:phyx-mini-0631:phyxminiproblems-problemphyxmini0631-massquantity, def:physics:phyx-mini-0631:phyxminiproblems-problemphyxmini0631-massinkilograms}
  A mass expressed as a dimensionless multiple of one atomic mass unit
\end{definition}
\begin{proof}
  This is the declared model definition of mass In Atomic Mass Units.
\end{proof}
\begin{definition}[moment Of Inertia In Kilogram Meters Squared]
  \label{def:physics:phyx-mini-0631:phyxminiproblems-problemphyxmini0631-momentofinertiainkilogrammeterssquared}
  \lean{PhyXMiniProblems.ProblemPhyXMini0631.momentOfInertiaInKilogramMetersSquared}
  \uses{def:physics:phyx-mini-0631:phyxminiproblems-problemphyxmini0631-momentofinertiaquantity}
  Coherent-SI readout of a moment of inertia, in kg m²
\end{definition}
\begin{proof}
  This is the declared model definition of moment Of Inertia In Kilogram Meters Squared.
\end{proof}
\begin{definition}[action In Joule Seconds]
  \label{def:physics:phyx-mini-0631:phyxminiproblems-problemphyxmini0631-actioninjouleseconds}
  \lean{PhyXMiniProblems.ProblemPhyXMini0631.actionInJouleSeconds}
  \uses{def:physics:phyx-mini-0631:phyxminiproblems-problemphyxmini0631-actionquantity}
  Coherent-SI readout of an action, in joule-seconds
\end{definition}
\begin{proof}
  This is the declared model definition of action In Joule Seconds.
\end{proof}
\begin{definition}[energy In Joules]
  \label{def:physics:phyx-mini-0631:phyxminiproblems-problemphyxmini0631-energyinjoules}
  \lean{PhyXMiniProblems.ProblemPhyXMini0631.energyInJoules}
  Coherent-SI readout of an energy, in joules
\end{definition}
\begin{proof}
  This is the declared model definition of energy In Joules.
\end{proof}
\begin{definition}[energy In Electron Volts]
  \label{def:physics:phyx-mini-0631:phyxminiproblems-problemphyxmini0631-energyinelectronvolts}
  \lean{PhyXMiniProblems.ProblemPhyXMini0631.energyInElectronVolts}
  \uses{def:physics:phyx-mini-0631:phyxminiproblems-problemphyxmini0631-energyinjoules}
  Electron-volt readout of a physical energy
\end{definition}
\begin{proof}
  This is the declared model definition of energy In Electron Volts.
\end{proof}
\begin{definition}[speed In Meters Per Second]
  \label{def:physics:phyx-mini-0631:phyxminiproblems-problemphyxmini0631-speedinmeterspersecond}
  \lean{PhyXMiniProblems.ProblemPhyXMini0631.speedInMetersPerSecond}
  \uses{def:physics:phyx-mini-0631:phyxminiproblems-problemphyxmini0631-speedquantity}
  Coherent-SI readout of a speed, in metres per second
\end{definition}
\begin{proof}
  This is the declared model definition of speed In Meters Per Second.
\end{proof}
\begin{definition}[Molecule Candidate]
  \label{def:physics:phyx-mini-0631:phyxminiproblems-problemphyxmini0631-moleculecandidate}
  \lean{PhyXMiniProblems.ProblemPhyXMini0631.MoleculeCandidate}
  The three possible gases stated in the problem
\end{definition}
\begin{proof}
  This is the declared model definition of Molecule Candidate.
\end{proof}
\begin{definition}[Chemical Element]
  \label{def:physics:phyx-mini-0631:phyxminiproblems-problemphyxmini0631-chemicalelement}
  \lean{PhyXMiniProblems.ProblemPhyXMini0631.ChemicalElement}
  Chemical-element labels visible beside the atoms in the figure
\end{definition}
\begin{proof}
  This is the declared model definition of Chemical Element.
\end{proof}
\begin{definition}[Atom Site]
  \label{def:physics:phyx-mini-0631:phyxminiproblems-problemphyxmini0631-atomsite}
  \lean{PhyXMiniProblems.ProblemPhyXMini0631.AtomSite}
  Left and right atomic sites of a depicted diatomic molecule
\end{definition}
\begin{proof}
  This is the declared model definition of Atom Site.
\end{proof}
\begin{definition}[Figure Row]
  \label{def:physics:phyx-mini-0631:phyxminiproblems-problemphyxmini0631-figurerow}
  \lean{PhyXMiniProblems.ProblemPhyXMini0631.FigureRow}
  Vertical positions of the three molecular diagrams in image 631.png
\end{definition}
\begin{proof}
  This is the declared model definition of Figure Row.
\end{proof}
\begin{definition}[Atom Color]
  \label{def:physics:phyx-mini-0631:phyxminiproblems-problemphyxmini0631-atomcolor}
  \lean{PhyXMiniProblems.ProblemPhyXMini0631.AtomColor}
  Atom colors used in the raster to distinguish the elements
\end{definition}
\begin{proof}
  This is the declared model definition of Atom Color.
\end{proof}
\begin{definition}[Illumination Kind]
  \label{def:physics:phyx-mini-0631:phyxminiproblems-problemphyxmini0631-illuminationkind}
  \lean{PhyXMiniProblems.ProblemPhyXMini0631.IlluminationKind}
  The illumination named by the physical scenario
\end{definition}
\begin{proof}
  This is the declared model definition of Illumination Kind.
\end{proof}
\begin{definition}[Spectrum Kind]
  \label{def:physics:phyx-mini-0631:phyxminiproblems-problemphyxmini0631-spectrumkind}
  \lean{PhyXMiniProblems.ProblemPhyXMini0631.SpectrumKind}
  The observed optical spectrum is an absorption rather than emission spectrum
\end{definition}
\begin{proof}
  This is the declared model definition of Spectrum Kind.
\end{proof}
\begin{definition}[Vibrational Level]
  \label{def:physics:phyx-mini-0631:phyxminiproblems-problemphyxmini0631-vibrationallevel}
  \lean{PhyXMiniProblems.ProblemPhyXMini0631.VibrationalLevel}
  Lower and upper vibrational manifolds of the observed band
\end{definition}
\begin{proof}
  This is the declared model definition of Vibrational Level.
\end{proof}
\begin{definition}[Absorption Band Label]
  \label{def:physics:phyx-mini-0631:phyxminiproblems-problemphyxmini0631-absorptionbandlabel}
  \lean{PhyXMiniProblems.ProblemPhyXMini0631.AbsorptionBandLabel}
  Labels for the five bands, named by their wavelengths in units of 10⁻⁴ μm. The declaration names are labels only; they do not encode any rotational assignment
\end{definition}
\begin{proof}
  This is the declared model definition of Absorption Band Label.
\end{proof}
\begin{definition}[Diatomic Candidates Figure]
  \label{def:physics:phyx-mini-0631:phyxminiproblems-problemphyxmini0631-diatomiccandidatesfigure}
  \lean{PhyXMiniProblems.ProblemPhyXMini0631.DiatomicCandidatesFigure}
  \uses{def:physics:phyx-mini-0631:phyxminiproblems-problemphyxmini0631-moleculecandidate, def:physics:phyx-mini-0631:phyxminiproblems-problemphyxmini0631-chemicalelement, def:physics:phyx-mini-0631:phyxminiproblems-problemphyxmini0631-atomsite, def:physics:phyx-mini-0631:phyxminiproblems-problemphyxmini0631-figurerow, def:physics:phyx-mini-0631:phyxminiproblems-problemphyxmini0631-atomcolor}
  Literal content carried by the primary molecular-candidate figure. Its mass numbers and bond-length texts are scalar raster labels; independent physical masses and lengths are stored in the setup below
\end{definition}
\begin{proof}
  This is the declared model definition of Diatomic Candidates Figure.
\end{proof}
\begin{definition}[Absorption Band]
  \label{def:physics:phyx-mini-0631:phyxminiproblems-problemphyxmini0631-absorptionband}
  \lean{PhyXMiniProblems.ProblemPhyXMini0631.AbsorptionBand}
  \uses{def:physics:phyx-mini-0631:phyxminiproblems-problemphyxmini0631-lengthquantity}
  Independent observables and quantum-number roles attached to one dark band. The energy is not defined from its wavelength, and the final quantum number is not defined from its label
\end{definition}
\begin{proof}
  This is the declared model definition of Absorption Band.
\end{proof}
\begin{definition}[Diatomic Absorption Setup]
  \label{def:physics:phyx-mini-0631:phyxminiproblems-problemphyxmini0631-diatomicabsorptionsetup}
  \lean{PhyXMiniProblems.ProblemPhyXMini0631.DiatomicAbsorptionSetup}
  \uses{def:physics:phyx-mini-0631:phyxminiproblems-problemphyxmini0631-lengthquantity, def:physics:phyx-mini-0631:phyxminiproblems-problemphyxmini0631-massquantity, def:physics:phyx-mini-0631:phyxminiproblems-problemphyxmini0631-momentofinertiaquantity, def:physics:phyx-mini-0631:phyxminiproblems-problemphyxmini0631-actionquantity, def:physics:phyx-mini-0631:phyxminiproblems-problemphyxmini0631-speedquantity, def:physics:phyx-mini-0631:phyxminiproblems-problemphyxmini0631-moleculecandidate, def:physics:phyx-mini-0631:phyxminiproblems-problemphyxmini0631-atomsite, def:physics:phyx-mini-0631:phyxminiproblems-problemphyxmini0631-illuminationkind, def:physics:phyx-mini-0631:phyxminiproblems-problemphyxmini0631-spectrumkind, def:physics:phyx-mini-0631:phyxminiproblems-problemphyxmini0631-vibrationallevel, def:physics:phyx-mini-0631:phyxminiproblems-problemphyxmini0631-absorptionbandlabel, def:physics:phyx-mini-0631:phyxminiproblems-problemphyxmini0631-diatomiccandidatesfigure, def:physics:phyx-mini-0631:phyxminiproblems-problemphyxmini0631-absorptionband}
  The gas candidates, their rigid-rotor parameters, and the measured spectrum. The vibrational term energies and full rovibrational energies are independent dimensionful fields related only by the governing laws below
\end{definition}
\begin{proof}
  This is the declared model definition of Diatomic Absorption Setup.
\end{proof}
\begin{definition}[Matches Gas Absorption Scenario]
  \label{def:physics:phyx-mini-0631:phyxminiproblems-problemphyxmini0631-matchesgasabsorptionscenario}
  \lean{PhyXMiniProblems.ProblemPhyXMini0631.MatchesGasAbsorptionScenario}
  \uses{def:physics:phyx-mini-0631:phyxminiproblems-problemphyxmini0631-diatomicabsorptionsetup}
  Qualitative information stated in the prose scenario
\end{definition}
\begin{proof}
  This is the declared model definition of Matches Gas Absorption Scenario.
\end{proof}
\begin{definition}[Matches Primary Diatomic Candidates Figure]
  \label{def:physics:phyx-mini-0631:phyxminiproblems-problemphyxmini0631-matchesprimarydiatomiccandidatesfigure}
  \lean{PhyXMiniProblems.ProblemPhyXMini0631.MatchesPrimaryDiatomicCandidatesFigure}
  \uses{def:physics:phyx-mini-0631:phyxminiproblems-problemphyxmini0631-lengthinpicometers, def:physics:phyx-mini-0631:phyxminiproblems-problemphyxmini0631-diatomicabsorptionsetup}
  Exact labels and qualitative geometry transcribed from the primary image. The final three fields connect each displayed bond-length readout to the corresponding independent physical candidate length. Nothing here concerns a spectral-line assignment or absorbed-energy answer
\end{definition}
\begin{proof}
  This is the declared model definition of Matches Primary Diatomic Candidates Figure.
\end{proof}
\begin{definition}[Agrees With Four Decimal Micrometer Readout]
  \label{def:physics:phyx-mini-0631:phyxminiproblems-problemphyxmini0631-agreeswithfourdecimalmicrometerreadout}
  \lean{PhyXMiniProblems.ProblemPhyXMini0631.AgreesWithFourDecimalMicrometerReadout}
  \uses{def:physics:phyx-mini-0631:phyxminiproblems-problemphyxmini0631-lengthquantity, def:physics:phyx-mini-0631:phyxminiproblems-problemphyxmini0631-lengthinmicrometers}
  Agreement with a wavelength printed to four decimal places in micrometres
\end{definition}
\begin{proof}
  This is the declared model definition of Agrees With Four Decimal Micrometer Readout.
\end{proof}
\begin{definition}[Matches Given Absorption Wavelengths]
  \label{def:physics:phyx-mini-0631:phyxminiproblems-problemphyxmini0631-matchesgivenabsorptionwavelengths}
  \lean{PhyXMiniProblems.ProblemPhyXMini0631.MatchesGivenAbsorptionWavelengths}
  \uses{def:physics:phyx-mini-0631:phyxminiproblems-problemphyxmini0631-diatomicabsorptionsetup, def:physics:phyx-mini-0631:phyxminiproblems-problemphyxmini0631-agreeswithfourdecimalmicrometerreadout}
  The five wavelength readouts supplied in the problem statement
\end{definition}
\begin{proof}
  This is the declared model definition of Matches Given Absorption Wavelengths.
\end{proof}
\begin{definition}[Uses Textbook Reference Data]
  \label{def:physics:phyx-mini-0631:phyxminiproblems-problemphyxmini0631-usestextbookreferencedata}
  \lean{PhyXMiniProblems.ProblemPhyXMini0631.UsesTextbookReferenceData}
  \uses{def:physics:phyx-mini-0631:phyxminiproblems-problemphyxmini0631-massinkilograms, def:physics:phyx-mini-0631:phyxminiproblems-problemphyxmini0631-massinatomicmassunits, def:physics:phyx-mini-0631:phyxminiproblems-problemphyxmini0631-actioninjouleseconds, def:physics:phyx-mini-0631:phyxminiproblems-problemphyxmini0631-speedinmeterspersecond, def:physics:phyx-mini-0631:phyxminiproblems-problemphyxmini0631-diatomicabsorptionsetup}
  Atomic-mass-number approximation and universal constants used by the textbook model. Physlib grounds ℏ, the electron volt, and the speed of light. This reference data does not select an absorbing gas, spectral line, or answer
\end{definition}
\begin{proof}
  This is the declared model definition of Uses Textbook Reference Data.
\end{proof}
\begin{definition}[Has Physical Diatomic Absorption Parameters]
  \label{def:physics:phyx-mini-0631:phyxminiproblems-problemphyxmini0631-hasphysicaldiatomicabsorptionparameters}
  \lean{PhyXMiniProblems.ProblemPhyXMini0631.HasPhysicalDiatomicAbsorptionParameters}
  \uses{def:physics:phyx-mini-0631:phyxminiproblems-problemphyxmini0631-lengthinmeters, def:physics:phyx-mini-0631:phyxminiproblems-problemphyxmini0631-massinkilograms, def:physics:phyx-mini-0631:phyxminiproblems-problemphyxmini0631-momentofinertiainkilogrammeterssquared, def:physics:phyx-mini-0631:phyxminiproblems-problemphyxmini0631-actioninjouleseconds, def:physics:phyx-mini-0631:phyxminiproblems-problemphyxmini0631-energyinjoules, def:physics:phyx-mini-0631:phyxminiproblems-problemphyxmini0631-speedinmeterspersecond, def:physics:phyx-mini-0631:phyxminiproblems-problemphyxmini0631-diatomicabsorptionsetup}
  Positivity and nondegeneracy conditions selecting physical parameters
\end{definition}
\begin{proof}
  This is the declared model definition of Has Physical Diatomic Absorption Parameters.
\end{proof}
\begin{definition}[Satisfies Diatomic Moment Of Inertia Laws]
  \label{def:physics:phyx-mini-0631:phyxminiproblems-problemphyxmini0631-satisfiesdiatomicmomentofinertialaws}
  \lean{PhyXMiniProblems.ProblemPhyXMini0631.SatisfiesDiatomicMomentOfInertiaLaws}
  \uses{def:physics:phyx-mini-0631:phyxminiproblems-problemphyxmini0631-lengthinmeters, def:physics:phyx-mini-0631:phyxminiproblems-problemphyxmini0631-massinkilograms, def:physics:phyx-mini-0631:phyxminiproblems-problemphyxmini0631-momentofinertiainkilogrammeterssquared, def:physics:phyx-mini-0631:phyxminiproblems-problemphyxmini0631-diatomicabsorptionsetup}
  Two-body center-of-mass geometry for every candidate: the reduced mass obeys μ(m₁+m₂)=m₁m₂, and a rigid diatomic of separation r has I=μr²
\end{definition}
\begin{proof}
  This is the declared model definition of Satisfies Diatomic Moment Of Inertia Laws.
\end{proof}
\begin{definition}[Satisfies Rigid Rotor Absorption Laws]
  \label{def:physics:phyx-mini-0631:phyxminiproblems-problemphyxmini0631-satisfiesrigidrotorabsorptionlaws}
  \lean{PhyXMiniProblems.ProblemPhyXMini0631.SatisfiesRigidRotorAbsorptionLaws}
  \uses{def:physics:phyx-mini-0631:phyxminiproblems-problemphyxmini0631-lengthinmeters, def:physics:phyx-mini-0631:phyxminiproblems-problemphyxmini0631-momentofinertiainkilogrammeterssquared, def:physics:phyx-mini-0631:phyxminiproblems-problemphyxmini0631-actioninjouleseconds, def:physics:phyx-mini-0631:phyxminiproblems-problemphyxmini0631-energyinjoules, def:physics:phyx-mini-0631:phyxminiproblems-problemphyxmini0631-speedinmeterspersecond, def:physics:phyx-mini-0631:phyxminiproblems-problemphyxmini0631-diatomicabsorptionsetup}
  The governing rigid-rotor and photon laws: within either vibrational manifold, E(v,l) = G(v) + ℏ² l(l+1)/(2I); a band photon supplies the upper-minus-lower rovibrational gap; Eγ λ = h c = 2πℏc; the dipole rotational selection rule is Δl = ±1; and five distinct lines originate in the lowest states l = 0,1,2. The final two clauses express the prose statement that the five lines are the transitions from the lowest rotational states. They do not associate any particular wavelength label with any particular transition
\end{definition}
\begin{proof}
  This is the declared model definition of Satisfies Rigid Rotor Absorption Laws.
\end{proof}
\begin{definition}[Answer Choice]
  \label{def:physics:phyx-mini-0631:phyxminiproblems-problemphyxmini0631-answerchoice}
  \lean{PhyXMiniProblems.ProblemPhyXMini0631.AnswerChoice}
  Labels of the four absorbed-energy choices in the source
\end{definition}
\begin{proof}
  This is the declared model definition of Answer Choice.
\end{proof}
\begin{definition}[displayed Energy Electron Volts]
  \label{def:physics:phyx-mini-0631:phyxminiproblems-problemphyxmini0631-displayedenergyelectronvolts}
  \lean{PhyXMiniProblems.ProblemPhyXMini0631.displayedEnergyElectronVolts}
  \uses{def:physics:phyx-mini-0631:phyxminiproblems-problemphyxmini0631-answerchoice}
  The energy printed beside each answer label, in electron volts
\end{definition}
\begin{proof}
  This is the declared model definition of displayed Energy Electron Volts.
\end{proof}
\begin{definition}[recorded Dataset Answer]
  \label{def:physics:phyx-mini-0631:phyxminiproblems-problemphyxmini0631-recordeddatasetanswer}
  \lean{PhyXMiniProblems.ProblemPhyXMini0631.recordedDatasetAnswer}
  \uses{def:physics:phyx-mini-0631:phyxminiproblems-problemphyxmini0631-answerchoice}
  Dataset metadata recording answer A; deliberately not used as a premise
\end{definition}
\begin{proof}
  This is the declared model definition of recorded Dataset Answer.
\end{proof}
\begin{definition}[answer Energy Error Electron Volts]
  \label{def:physics:phyx-mini-0631:phyxminiproblems-problemphyxmini0631-answerenergyerrorelectronvolts}
  \lean{PhyXMiniProblems.ProblemPhyXMini0631.answerEnergyErrorElectronVolts}
  \uses{def:physics:phyx-mini-0631:phyxminiproblems-problemphyxmini0631-energyinelectronvolts, def:physics:phyx-mini-0631:phyxminiproblems-problemphyxmini0631-absorptionbandlabel, def:physics:phyx-mini-0631:phyxminiproblems-problemphyxmini0631-diatomicabsorptionsetup, def:physics:phyx-mini-0631:phyxminiproblems-problemphyxmini0631-answerchoice, def:physics:phyx-mini-0631:phyxminiproblems-problemphyxmini0631-displayedenergyelectronvolts}
  Error between a band's physical energy and a displayed choice
\end{definition}
\begin{proof}
  This is the declared model definition of answer Energy Error Electron Volts.
\end{proof}
\begin{definition}[Agrees With Displayed Energy]
  \label{def:physics:phyx-mini-0631:phyxminiproblems-problemphyxmini0631-agreeswithdisplayedenergy}
  \lean{PhyXMiniProblems.ProblemPhyXMini0631.AgreesWithDisplayedEnergy}
  \uses{def:physics:phyx-mini-0631:phyxminiproblems-problemphyxmini0631-absorptionbandlabel, def:physics:phyx-mini-0631:phyxminiproblems-problemphyxmini0631-diatomicabsorptionsetup, def:physics:phyx-mini-0631:phyxminiproblems-problemphyxmini0631-answerchoice, def:physics:phyx-mini-0631:phyxminiproblems-problemphyxmini0631-answerenergyerrorelectronvolts}
  Agreement with the source answer to 10⁻⁵ eV. This tolerance covers the four-decimal-place wavelength readout and the rounding conventions used in the answer list
\end{definition}
\begin{proof}
  This is the declared model definition of Agrees With Displayed Energy.
\end{proof}
\begin{definition}[Is Unique Closest Energy Choice]
  \label{def:physics:phyx-mini-0631:phyxminiproblems-problemphyxmini0631-isuniqueclosestenergychoice}
  \lean{PhyXMiniProblems.ProblemPhyXMini0631.IsUniqueClosestEnergyChoice}
  \uses{def:physics:phyx-mini-0631:phyxminiproblems-problemphyxmini0631-absorptionbandlabel, def:physics:phyx-mini-0631:phyxminiproblems-problemphyxmini0631-diatomicabsorptionsetup, def:physics:phyx-mini-0631:phyxminiproblems-problemphyxmini0631-answerchoice, def:physics:phyx-mini-0631:phyxminiproblems-problemphyxmini0631-answerenergyerrorelectronvolts}
  The selected energy is strictly closer than every other displayed choice
\end{definition}
\begin{proof}
  This is the declared model definition of Is Unique Closest Energy Choice.
\end{proof}
\begin{lemma}[band42972 is transition one to zero]
  \label{lem:physics:phyx-mini-0631:phyxminiproblems-problemphyxmini0631-band42972-is-transition-one-to-zero}
  \lean{PhyXMiniProblems.ProblemPhyXMini0631.band42972_is_transition_one_to_zero}
  \uses{def:physics:phyx-mini-0631:phyxminiproblems-problemphyxmini0631-diatomicabsorptionsetup, def:physics:phyx-mini-0631:phyxminiproblems-problemphyxmini0631-matchesgivenabsorptionwavelengths, def:physics:phyx-mini-0631:phyxminiproblems-problemphyxmini0631-hasphysicaldiatomicabsorptionparameters, def:physics:phyx-mini-0631:phyxminiproblems-problemphyxmini0631-satisfiesrigidrotorabsorptionlaws}
  The ordered wavelengths, positive rigid-rotor spacing, Δl = ±1, and complete set of distinct lowest-state transitions identify the first line on the long-wavelength (P) side as l = 1 → 0. This is a derived assignment, not a premise
\end{lemma}
\begin{proof}
  The conclusion follows from the governing relations and figure readouts named in the statement of band42972 is transition one to zero.
\end{proof}
\begin{lemma}[band42972 energy is choice A]
  \label{lem:physics:phyx-mini-0631:phyxminiproblems-problemphyxmini0631-band42972-energy-is-choicea}
  \lean{PhyXMiniProblems.ProblemPhyXMini0631.band42972_energy_is_choiceA}
  \uses{def:physics:phyx-mini-0631:phyxminiproblems-problemphyxmini0631-diatomicabsorptionsetup, def:physics:phyx-mini-0631:phyxminiproblems-problemphyxmini0631-matchesgivenabsorptionwavelengths, def:physics:phyx-mini-0631:phyxminiproblems-problemphyxmini0631-usestextbookreferencedata, def:physics:phyx-mini-0631:phyxminiproblems-problemphyxmini0631-hasphysicaldiatomicabsorptionparameters, def:physics:phyx-mini-0631:phyxminiproblems-problemphyxmini0631-satisfiesrigidrotorabsorptionlaws, def:physics:phyx-mini-0631:phyxminiproblems-problemphyxmini0631-agreeswithdisplayedenergy, def:physics:phyx-mini-0631:phyxminiproblems-problemphyxmini0631-isuniqueclosestenergychoice}
  Applying Eγλ = hc to the measured 4.2972 μm wavelength gives an energy within 10⁻⁵ eV of 0.288528 eV, and choice A is uniquely closest
\end{lemma}
\begin{proof}
  The conclusion follows from the governing relations and figure readouts named in the statement of band42972 energy is choice A.
\end{proof}
% --- Archon declaration topology end ---
% --- Archon physics formalization source end ---
